\documentclass[11pt]{article}

\usepackage[margin=1in]{geometry}
\usepackage{amsmath,amssymb,amsthm}
\usepackage{booktabs}
\usepackage{enumitem}
\usepackage{hyperref}
\usepackage{microtype}

\hypersetup{colorlinks=true,linkcolor=blue!60!black,citecolor=blue!60!black,urlcolor=blue!60!black}

\newtheorem{theorem}{Theorem}
\newtheorem{lemma}[theorem]{Lemma}
\newtheorem{proposition}[theorem]{Proposition}
\newtheorem{corollary}[theorem]{Corollary}
\theoremstyle{definition}
\newtheorem{definition}[theorem]{Definition}
\newtheorem{openproblem}{Open Problem}
\theoremstyle{remark}
\newtheorem*{remark}{Remark}

\newcommand{\Z}{\mathbb{Z}}
\newcommand{\F}{\mathbb{F}}
\newcommand{\GL}{\mathrm{GL}}
\newcommand{\im}{\mathrm{im}}
\newcommand{\rank}{\mathrm{rank}}
\newcommand{\unrank}{\mathrm{unrank}}

\title{\textbf{Ranked Digit Layouts}\\[4pt]
\large An Algebra for Tensor Memory Mapping\\
with Arbitrary Dimension Sizes}
\author{}
\date{}

\begin{document}
\maketitle
\thispagestyle{empty}

\begin{abstract}
Linear Layouts over $\F_2$ give tensor compilers closed-form resolution of memory mapping, vectorization, and hardware lowering---but only when every dimension is a power of two.
We define \emph{Ranked Digit Layouts}: factor every dimension into primes, operate on the resulting digit stream with a small set of typed transforms, and decode.
The framework handles arbitrary dimension sizes exactly, composes cleanly, and admits $O(1)$-in-volume algorithms for contiguity, intersection, and kernel analysis.
We give a constructive algorithm for provably optimal bank-conflict-free swizzle synthesis within the algebra.
\end{abstract}

%% ========================================================================
\part*{I\quad The Problem}
%% ========================================================================

\section{Setting}

A tensor compiler must map three coordinate spaces onto one another: logical tensors, physical memory, and hardware topology.
Over $\F_2$, every coordinate is a bit vector, every layout is a binary matrix, and composition, inversion, intersection, and contiguity analysis all reduce to linear algebra over $\F_2$ in time independent of the tensor volume.

The cost: every logical dimension must be padded to a power of two.
A $48 \times 48$ tile becomes $64 \times 64$, wasting $56\%$ of memory.

\begin{openproblem}
Does there exist an algebraic layout representation that preserves the analytical tractability of $\F_2$ linear layouts while natively supporting arbitrary dimension sizes?
\end{openproblem}

\noindent We ask for nine properties, grouped into three classes.

\medskip
\noindent\textbf{I.\;Structural.}\;
\textbf{(1)}~Domain/codomain agnosticism: no padding, no masking.
\textbf{(2)}~Universal composition: layouts compose within the algebra.
\textbf{(3)}~Shape closure: transpose, reshape, and broadcast are closed operations on layouts.

\medskip
\noindent\textbf{II.\;Analytical tractability}\,(all $O(1)$ in the tensor volume).
\textbf{(4)}~Tractable intersections for data routing.
\textbf{(5)}~Tractable kernels for broadcast detection.
\textbf{(6)}~Tractable contiguity for vectorization.

\medskip
\noindent\textbf{III.\;Constructive synthesis.}\;
\textbf{(7)}~Intrinsic factorization: algebraic division by a hardware primitive.
\textbf{(8)}~Optimal swizzle synthesis with automorphic closure.
\textbf{(9)}~Topology-aware cost resolution.

\medskip
\noindent Solving this bridges the geometric flexibility of the Polyhedral model and the microscopic efficiency of Linear Layouts over~$\F_2$.

%% ========================================================================
\part*{II\quad The Construction}
%% ========================================================================

\section{Spaces and Digit Streams}

\begin{definition}[Space]
A \emph{space} is a finite product of cyclic groups
$A = \Z_{n_0} \times \Z_{n_1} \times \cdots \times \Z_{n_{k-1}}$
with $|A| = \prod_i n_i$.
The bounds $n_i$ are arbitrary positive integers.
\end{definition}

\begin{definition}[Digit stream]
Factor each $n_i$ into primes: $n_i = p_{i,0}\, p_{i,1}\cdots p_{i,\,r_i-1}$.
A coordinate $c_i \in \Z_{n_i}$ is encoded as digits $(d_0,\ldots,d_{r_i-1})$ with $d_j \in \Z_{p_{i,j}}$ via the standard little-endian mixed-radix representation:
\[
c_i = d_0 + p_{i,0}\,d_1 + p_{i,0}\,p_{i,1}\,d_2 + \cdots
\]
Concatenating the digit vectors of all axes gives a flat \emph{digit stream} of length $R = \sum_i r_i$.
The map $\rank: A \to \Z_{p_0}\times\cdots\times\Z_{p_{R-1}}$ is a bijection; its inverse is $\unrank$.
\end{definition}

\begin{remark}
Every digit lives in a prime field $\F_p$.
There are no carries between digit positions, because each position is reduced independently modulo its own prime.
This is the key difference from working over $\Z_N$ directly.
\end{remark}

\section{Transforms}

A \emph{transform} maps one digit stream to another.
Six primitives suffice.

\begin{definition}[Identity / Refactor]
$T(\mathbf{d}) = \mathbf{d}$.
A \emph{Refactor} changes only how digits are grouped into axes; the stream itself is unchanged.
Invertible; kernel size~$1$.
\end{definition}

\begin{definition}[Permute]
A bijection defined by a permutation $\sigma$: $T(\mathbf{d})_i = d_{\sigma(i)}$.
Axis transposition is a special case.
Invertible; kernel size~$1$.
\end{definition}

\begin{definition}[Project]
Drops digit positions outside a keep-set $K$: $T(\mathbf{d})_i = d_{K_i}$.
Non-injective; kernel size $\prod_{i\notin K} p_i$.
Models broadcasting and slicing.
\end{definition}

\begin{definition}[Embed]
Places source digits at specified positions in a larger stream, filling the rest with constants.
Right inverse of Project.
\end{definition}

\begin{definition}[BlockAffine]
Selects digit positions $P$ sharing a common prime radix $p$ and applies an affine map:
\[
T(\mathbf{d})_{P_i} = \Bigl(\sum_j M_{ij}\,d_{P_j} + c_i\Bigr) \bmod p,
\qquad M \in \GL(|P|, \F_p).
\]
Digits outside $P$ pass through.
XOR swizzle is the case $p=2$, $c=0$.
Invertible; kernel size~$1$.
\end{definition}

\begin{definition}[Compose]
Sequential application: $(T_2 \circ T_1)(\mathbf{d}) = T_2(T_1(\mathbf{d}))$.
Requires output radices of $T_1$ to match input radices of $T_2$.
Associative; forms a monoid with Identity.
\end{definition}

\section{Ranked Digit Layouts}

\begin{definition}
A \emph{ranked digit layout} $L: A \to B$ is a triple $(A,\,T,\,B)$ where $T$ is a transform whose input radices match the digit stream of $A$ and whose output radices match the digit stream of $B$.
The induced coordinate map is
$L(c) = \unrank_B(T(\rank_A(c)))$.
\end{definition}

When $B = \Z_N$ is flat, $L$ gives a memory address.
When $A$ is a hardware thread space and $B$ is a tile, $L$ describes a thread assignment.

\section{Composition and Inversion}

Given $L_1: A \to B$ and $L_2: B \to C$, the composition $L_2 \circ L_1: A \to C$ is defined by composing transforms.
Inversion of $L$ exists when $T$ is invertible.

\emph{Algebraic division.}
Given the composed layout $L \circ H$ and an invertible hardware map~$H$, the logical layout is recovered as $(L \circ H) \circ H^{-1} = L$.
This is the mechanism for hardware lowering: given a hardware primitive, divide it out.

\section{Analytical Contiguity}

\begin{proposition}
For $L: A \to \Z_N$ and a logical axis $a$, the contiguity along $a$---the largest $w$ such that stepping through the first $w$ elements along $a$ from the origin produces consecutive flat addresses---is computable in $O(\Omega(n_a))$ transform evaluations.
\end{proposition}

\begin{proof}
Walk the digits of axis $a$ from least to most significant.
At digit $d$ with radix $r_d$:
\begin{enumerate}[nosep]
\item Evaluate the flat address with only digit $d$ set to $1$; call the stride $s_d$.
\item If $s_d$ does not equal the running expected stride, stop.
\item Verify linearity: the address with digit $d$ set to $r_d - 1$ must equal $(r_d - 1)\,s_d$ above the base.
\item Multiply the contiguity width by $r_d$ and update the expected stride.
\end{enumerate}
Non-linearity at digit $d$ implies non-contiguity at all higher digits, because higher digits control strides that are multiples of the current position weight.
Hence the procedure returns the \emph{largest} contiguous width.
The number of steps is $\Omega(n_a)$, the number of prime factors of the axis size---logarithmic in~$n_a$.
\end{proof}

\section{Swizzle Synthesis}

Let $L: A \to V$ be a layout into physical memory $V = \Z_N$, with $B$ memory banks indexed by $a \bmod B$ where $B \mid N$.
A warp of $W$ threads accesses $V$ simultaneously.
We seek a bijective automorphism $S: V \to V$ such that $S \circ L$ is conflict-free and remains in the algebra.

\subsection*{Setup}

Factor $B$ into primes.
For each prime $p$, let $n_p = v_p(N)$ (the number of radix-$p$ digits in $V$) and $b_p = v_p(B)$.
The bank index, read in digit space, is a projection $E_p: \F_p^{n_p} \to \F_p^{b_p}$ keeping the $b_p$ least-significant radix-$p$ digits.

The composed map $L \circ H$ (where $H$ is the hardware thread assignment) gives each thread a flat address whose radix-$p$ component is affine in the thread's radix-$p$ digits.
Let $\varphi_p: \F_p^{w_p} \to \F_p^{n_p}$ be the linear part.
Define the \emph{thread subspace} $\Delta_p = \im(\varphi_p)$.

Two threads conflict in the $p$-component iff $E_p(\varphi_p(t_1) - \varphi_p(t_2)) = 0$, i.e.\ their address difference lies in $\ker(E_p) \cap \Delta_p$.

\subsection*{The Result}

\begin{theorem}[Optimal swizzle bound]
\label{thm:swizzle}
A BlockAffine swizzle $S$ with matrix $M_p \in \GL(n_p, \F_p)$ yields a conflict factor for prime $p$ of
\[
C_p \;=\; p^{\;\dim(\ker(E_p \cdot M_p)\,\cap\,\Delta_p)}.
\]
The minimum achievable conflict factor is
\[
C_p^* \;=\; p^{\;\max(0,\; d_p - b_p)},
\qquad d_p = \dim(\Delta_p).
\]
The total conflict factor across all primes is $C = \prod_p C_p$.
\end{theorem}

\begin{proof}
The map $E_p \cdot M_p$ restricted to $\Delta_p$ is a linear map from a $d_p$-dimensional space to a $b_p$-dimensional space.
Its kernel has dimension at least $\max(0,\, d_p - b_p)$ by the rank--nullity theorem.
Hence $C_p \ge p^{\max(0,\,d_p - b_p)}$.
The construction below achieves this bound.
\end{proof}

\subsection*{Construction}

For each prime $p$ independently:

\begin{enumerate}[nosep]
\item Compute a basis $\{v_1,\ldots,v_{d_p}\}$ of $\Delta_p$.
\item Let $k = \min(d_p,\, b_p)$.
Pick $k$ linearly independent targets $u_1,\ldots,u_k \in \F_p^{b_p}$.
\item For each $u_i$, find a preimage $w_i$ with $E_p\,w_i = u_i$ (possible since $E_p$ is surjective).
\item Define $M_p$ by $M_p\,v_i = w_i$ for $i = 1,\ldots,k$.
Extend to a basis of $\F_p^{n_p}$ and complete $M_p$ to a full-rank matrix.
\end{enumerate}

By construction, $E_p \cdot M_p$ maps $\{v_1,\ldots,v_k\}$ to independent vectors, so $\ker(E_p \cdot M_p) \cap \Delta_p$ has the minimal possible dimension.
The cost is $O(n_p^3)$ field operations, with $n_p = O(\log_p N)$---independent of the tensor volume.

\begin{corollary}[GPU XOR swizzle recovery]
For $B = 32 = 2^5$, $W = 32 = 2^5$: only $p = 2$ is relevant, $b_2 = 5$, and generically $d_2 = 5$.
Since $d_2 = b_2$, we get $C_2^* = 1$: a zero-conflict swizzle exists.
The matrix $M_2$ is a binary invertible matrix---XOR swizzling.
\end{corollary}

\subsection*{Closure}

The swizzle $S$ is a BlockAffine automorphism on $V$'s digit stream.
Composing with $L$ appends a BlockAffine to the transform chain, remaining in the algebra.
Contiguity, kernel, and image analyses apply unchanged to $S \circ L$.

%% ========================================================================
\part*{III\quad Assessment}
%% ========================================================================

\section{Property Satisfaction}

\begin{table}[h]
\centering
\renewcommand{\arraystretch}{1.25}
\begin{tabular}{@{}clcc@{}}
\toprule
\textbf{\#} & \textbf{Property} & \textbf{Status} & \textbf{Section} \\
\midrule
1 & Domain/codomain agnosticism & \checkmark & 2 \\
2 & Universal composition & \checkmark & 4 \\
3 & Shape closure & \checkmark & 3 \\
4 & Tractable intersections & \checkmark & 3, 5 \\
5 & Tractable kernels & \checkmark & 3 \\
6 & Tractable contiguity & \checkmark & 5 \\
7 & Intrinsic factorization & \checkmark & 4 \\
8 & Swizzle synthesis \& closure & \checkmark & 6 \\
9 & Topology-aware cost & \textit{open} & --- \\
\bottomrule
\end{tabular}
\end{table}

\section{An Example: Reshape Without Padding}

Reshape $[6, 10] \to [12, 5]$, row-major.

\medskip
\noindent\emph{Source.}\;
Axis $0$: $6 = 2 \cdot 3$, digits $[\Z_2, \Z_3]$.\;
Axis $1$: $10 = 2 \cdot 5$, digits $[\Z_2, \Z_5]$.\\
Stream: $(\Z_2, \Z_3, \Z_2, \Z_5)$.

\medskip
\noindent\emph{Target.}\;
Axis $0$: $12 = 2 \cdot 2 \cdot 3$, digits $[\Z_2, \Z_2, \Z_3]$.\;
Axis $1$: $5$, digits $[\Z_5]$.\\
Stream: $(\Z_2, \Z_2, \Z_3, \Z_5)$.

\medskip
\noindent Under row-major ordering, the flat index is preserved.
The reshape is a Permute (moving the $\Z_2$ digit from axis~1 into axis~0) composed with a Refactor (regrouping axis boundaries).
The digit stream is rearranged; no ring structure changes; no padding is introduced.

\section{What Remains Open}

\textbf{Topology-aware cost resolution.}
When routing between spaces admits multiple valid solutions, the algebra does not yet provide a mechanism to minimize a hardware-topology-aware cost function.
This requires modeling physical spaces as hierarchical metric spaces and integrating cost-weighted selection into the underdetermined solve.

\medskip
\noindent\textbf{Canonicalization.}
Two composed transforms may be extensionally equal without being syntactically equal.
A polynomial-time canonical form for arbitrary transform chains has not been established.

\medskip
\noindent\textbf{Large primes.}
A dimension of prime size $p$ produces a single digit in $\F_p$.
The automorphism group is thin ($p-1$ scalar multiplications), limiting swizzle options.
In practice, ML tensor shapes are almost always highly composite, so this may be acceptable.

\bigskip
\hrule
\bigskip

\noindent When all dimension sizes are powers of two, every digit has radix~$2$, BlockAffine reduces to binary matrix--vector multiply, and the framework collapses exactly to Linear Layouts over~$\F_2$.
Ranked Digit Layouts are the natural generalization.

\end{document}
